\documentclass{article}
\usepackage[utf8]{inputenc}
\usepackage[T1]{fontenc}

\usepackage{amsmath,amssymb}
\usepackage{graphicx}
\usepackage{booktabs}
\usepackage{geometry}
\usepackage{enumerate}

\begin{document}


学术论文示例

摘要

本文研究了基于深度学习的图像分类方法。我们提出了一种新的卷积神经网络架构，在多个基准数据集上取得了最先进的性能。

\textbf{关键词: }深度学习, 图像分类, 卷积神经网络, 计算机视觉

1. 引言

近年来，深度学习技术在计算机视觉领域取得了显著进展。卷积神经网络 (CNN) 已成为图像分类的主流方法。

2. 方法

2.1 网络架构

我们提出的网络包含以下主要组件：

卷积层：用于特征提取

池化层：用于降维

全连接层：用于分类

2.2 损失函数

交叉熵损失: L = -∑(y log(ŷ)))

3. 实验结果

\begin{tabular}{|c|c|c|}
\hline
方法 & 准确率 & 参数量
\hline
ResNet-50 & 76.5\textbackslash{}% & 25.6M
\hline
VGG-16 & 72.8\textbackslash{}% & 138M
\hline
我们的方法 & 78.9\textbackslash{}% & 12.3M
\hline
\end{tabular}

4. 结论

本文提出了一种高效的图像分类方法。实验结果表明，我们的方法在准确率和效率上都优于现有方法。

参考文献

[1] He, K. et al. (2016). Deep Residual Learning for Image Recognition.

[2] Simonyan, K. & Zisserman, A. (2014). Very Deep Convolutional Networks.

\end{document}